% Chapitre « Validation » (axe 3). Sources : v2/docs/lab/reports/2026-11-14-axe3-*.md,
% v2/papers/axe3-validation/main.tex, v2/validation/preregistered/ (lecture seule),
% v2/docs/lab/decisions.md, v2/docs/lab/open.md, historique git. État relu le 2026-11-15.
\section{Validation}\label{sec:validation-ref}\label{val:sec}
\begingroup\setlength{\emergencystretch}{3em}% local à ce chapitre

Ce chapitre décrit comment le jumeau a été jugé, ce que les douze protocoles de l'axe~3 ont
établi, et ce qu'ils n'établissent pas. Le détail de chaque protocole est dans les rapports
\texttt{v2/docs/lab/reports/2026-11-14-axe3-n*.md}, abrégés ici \texttt{R/n1} à \texttt{R/n12}.
Le journal des décisions \texttt{v2/docs/lab/decisions.md} est abrégé \texttt{D}, l'article
\texttt{v2/papers/axe3-validation/main.tex} est abrégé \texttt{A}, et le registre des questions
ouvertes \texttt{v2/docs/lab/open.md} est abrégé \texttt{O}. Les numéros de ligne renvoient à
l'état du 2026-11-15.

\paragraph{Convention.} Chaque nombre porte sa provenance : \mesure{}~mesuré (par le projet ou
par la source citée), \derive{}~dérivé, \suppose{}~supposé (un choix, un seuil). Les années, les
identifiants (N1, K3, empreintes, commits) et les dates de calendrier ne sont pas étiquetés. Les
heures de commit sont lues dans git (\mesure). Les heures de gel « déclarées » sont celles que
donnent les rapports ; elles n'ont pas de preuve indépendante.

\subsection{Démarche}\label{val:demarche}

\paragraph{Pourquoi pré-enregistrer.} Le premier backtest du projet (K1, version~1) a retenu la
meilleure de trois variantes après avoir vu les résultats, pour un écart de \num{0.1}\mesure{}
point de MAPE (\texttt{A} l.58--62). C'est un biais de sélection. L'axe~3 inverse l'ordre : pour
chaque niveau de validation, un protocole écrit l'hypothèse, la statistique, le seuil et le test
qui le falsifierait, puis il est gelé avant tout calcul qui le concerne
(\texttt{v2/docs/lab/reports/2026-11-14-axe3-plan-validation.md}).

\paragraph{Gel par empreinte.} Une copie de chaque protocole est placée dans
\texttt{v2/\allowbreak validation/\allowbreak preregistered/}, et son empreinte SHA-256 est inscrite. Cette copie ne se
modifie plus ; la section « Résultat » est ajoutée au rapport, pas à la copie gelée. Le
2026-11-15, les \num{13}\mesure{} empreintes (plan et N1 à N12) ont été recalculées : toutes
concordent avec les valeurs inscrites (section~\ref{val:integrite}). Les copies gelées ne
diffèrent des rapports que par des chemins de liens relatifs, l'en-tête de N3 et la section
« Résultat » ajoutée ensuite (comparaison de \texttt{v2/\allowbreak validation/\allowbreak preregistered/} et de \texttt{v2/\allowbreak docs/\allowbreak lab/\allowbreak reports/}).

\paragraph{Baselines naïves.} Toute affirmation de précision se mesure contre la meilleure de deux
références naïves, la persistance (dernier effectif observé) et la tendance linéaire. C'est une
règle du projet (\texttt{CLAUDE.md}, « Accuracy claims need a backtest against a naive
baseline »). Pour N1, la statistique est
$\Delta = \mathrm{MAPE}(\text{candidat}) - \mathrm{MAPE}(\text{meilleure naïve})$, et le
candidat gagne si la borne haute de l'IC à 95~\% est négative (\texttt{R/n1} l.94--106). Pour les
prévisions probabilistes (N6, N7-C, N8), la statistique est l'écart de CRPS sur $\log N$
\citep{gneiting2007jasa}, et la calibration se lit par la couverture et le PIT
\citep{gneiting2007jrssb}. La SBC \citep{talts2018sbc} contrôle l'inférence dans des mondes
simulés.

\paragraph{Niveaux.} Les douze protocoles couvrent la prévision à 5\suppose{}~ans (N1 ; $h = 1$ à 10 en secondaire),
la cohérence à 2050 (N2), l'effet d'un levier (N3, N7-A), l'échelle communale (N4, N10 à N12),
des mondes simulés à vérité connue (N5), la prévision probabiliste (N6, N7-C, N8) et le signe des
boucles de rétroaction (N9). Les tableaux~\ref{val:tab-protocoles-a} et~\ref{val:tab-protocoles-b}
les résument. La colonne « Gel » donne l'heure déclarée, puis l'heure du commit qui contient la
copie gelée.

\begin{table}[p]
\centering
\footnotesize
\setlength{\tabcolsep}{3pt}
\renewcommand{\arraystretch}{1.15}
\begin{tabular}{@{}>{\raggedright\arraybackslash}p{0.5cm}>{\raggedright\arraybackslash}p{2.9cm}>{\raggedright\arraybackslash}p{3.3cm}>{\raggedright\arraybackslash}p{2.0cm}>{\raggedright\arraybackslash}p{3.6cm}>{\raggedright\arraybackslash}p{2.1cm}@{}}
\toprule
& \textbf{Hypothèse} & \textbf{Statistique et seuil} & \textbf{Gel} & \textbf{Résultat} & \textbf{Verdict} \\
\midrule
N1 & La variante par défaut prévoit l'effectif départemental mieux que la meilleure naïve &
$\Delta$ MAPE à $h = 5$\suppose, dép.\ avec $N \ge 30$\suppose ; bootstrap de
\num{2000}\suppose{} tirages ; gagne si borne haute $< 0$ &
02:54 décl. ; 03:25:50\mesure{} (316a93d), avant le calcul &
MAPE \num{8.68}\mesure{}~\% contre \num{9.30}\mesure{} (tendance) ; $\Delta = \num{-0.63}$
$[\num{-0.97} ; \num{-0.27}]$\mesure. Par origine : \num{+0.10}\mesure{} (2016),
\num{-0.22}\mesure{} (2018), \num{-1.22}\mesure{} (2021) &
\textbf{passé} au poolé, classement instable ; dépôt \textbf{en attente des données 2027} \\
\addlinespace
N2 & (A) reproduire DD76 en ratios $N(t)/N(2020)$ ; (B) cas limites L1--L3 &
A : $|\text{écart}| \le 2$\suppose{}~\% en 2030, $\le 3$\suppose{}~\% en 2040 et 2050 ; B : écart
$\le 10^{-9}$\suppose{} (L1, L3), $\le 1$\suppose{}~\% (L2) &
02:54 décl. ; 03:25:50\mesure, avant A (04:46\mesure) et B (08:08\mesure) &
A, forme affaiblie : 2050 CD \num{-14.1}\mesure, PH \num{-15.0}\mesure, SF
\num{-17.4}\mesure{}~\% ; B : L1 écart 0\mesure, L2 \num{-0.05} à \num{-0.46}\mesure{}~\% ;
ER~1319 : \num{-25.2}\mesure{}~\% en 2050 &
A \textbf{échoué} (forme fidèle \textbf{non jugeable}) ; B \textbf{passé} ; aucun chiffre 2050
citable \\
\addlinespace
N3 & E1 : la régulation 2009/2012 en zone~5 freine les IDE libéraux ($\beta < 0$) ; E2 : pente
des kinés négative en marche, remontée à l'arrêt &
T1 Hotelling $p \ge \num{0.10}$\suppose{} ; T2 placebo de date, IC $\ni 0$ ; T3--T4 placebos ;
T5 IC $< 0$ ; $\theta$ interprété seulement si T1--T4 passent &
02:54 décl. ; copie et résultats \textbf{dans le même commit} 316a93d &
E1 : T1 $p = \num{0.043}$\mesure{} ; effet \num{-30.3} $[\num{-34.1} ; \num{-26.4}]$\mesure{}
log-points ; $\theta = \num{0.80}$ $[\num{0.74} ; \num{0.87}]$\mesure. E2 : écart des pentes
\num{+3.1} $[\num{+1.1} ; \num{+5.1}]$\mesure{} pts/an &
E1 \textbf{échoué} (levier IDE non validé) ; E2 \textbf{passé} au dép. ; $\theta$ \textbf{non
jugeable} \\
\addlinespace
N4 & L'erreur communale à 5~ans décroît avec la taille ; $N_{\min} \ge 16$\derive &
$N_{\min}$ : P80 de $|$erreur$| \le 25$\suppose{}~\% ; bande P10--P90 couvrant
$[75 ; 85]$\suppose{}~\% \emph{dans chaque classe} &
02:54 décl. ; 03:25:50\mesure, avant le résultat (08:08\mesure) &
$N_{\min} = 51$\mesure{} (21\mesure{} sans les kinés) ; bande calibrée dans 14\mesure{} classes sur
21\mesure{} ; une inversion de monotonie\mesure{} &
$N_{\min}$ \textbf{passé} ; bande \textbf{échoué} ; sous oracle seulement \\
\addlinespace
N5 & Monde conforme (M1, M3) : variante exacte ; monde non conforme (M4--M7) : signe du biais
prévisible &
M1, M3 : $|$erreur$| \le 1$\suppose{}~\% à $h = 15$\suppose{} ; signe dans $\ge 16$\suppose{}
réplications sur 20\suppose{} ; M5 à 3\suppose{}~points de \num{-12.4}\suppose{}~\% &
03:46 décl. ; 03:55:53\mesure{} (72812c8), avant les résultats &
M1 max $4\cdot10^{-14}$\mesure{}~\% ; M3 9\mesure/20 sous 1~\% ; M4 \num{-34.5}\mesure{}~\%,
0\mesure/20 ; M5 \num{-19.8}\mesure{}~\% ; M6, M7 20\mesure/20 &
\textbf{échoué} (M3, M4, ampleur de M5) \\
\addlinespace
N6 & Bandes calibrées ; le candidat bat la naïve en CRPS &
CRPS sur $\log N$, critère de N1 ; couverture à 80~\% dans $[75 ; 85]$\suppose{}~\% ; SBC,
$p \ge \num{0.01}$\suppose{} (Bonferroni) &
03:53 décl. ; 03:55:53\mesure, avant les résultats &
279\derive{} couples (origine 2021) ; variante par défaut : $\Delta$CRPS $= \num{+0.0098}$
$[\num{+0.0037} ; \num{+0.0161}]$\mesure{} (\texttt{R/n6}), \num{+0.0109} $[\num{+0.0054} ;
\num{+0.0159}]$\mesure{} (\texttt{coh\_verdicts.json}) ; couverture « 80~\% » \num{56.6}\mesure{}~\% (de
\num{55.6} à \num{70.3}\mesure{}~\% pour les autres estimateurs) ; SBC entrées $p = 3\cdot10^{-46}$\mesure &
\textbf{échoué} ; SBC de la rétention passée \\
\bottomrule
\end{tabular}
\caption{Protocoles N1 à N6 : hypothèse, critère fixé d'avance, gel, résultat et verdict. Sources :
\texttt{R/n1} l.94--166, \texttt{R/n2} l.55--188, \texttt{R/n3} l.107--224, \texttt{R/n4}
l.71--170, \texttt{R/n5} l.42--113, \texttt{R/n6} l.47--157 ; heures de commit lues par
\texttt{git log -- v2/validation/preregistered}.}
\label{val:tab-protocoles-a}
\end{table}

\begin{table}[p]
\centering
\footnotesize
\setlength{\tabcolsep}{3pt}
\renewcommand{\arraystretch}{1.15}
\begin{tabular}{@{}>{\raggedright\arraybackslash}p{0.5cm}>{\raggedright\arraybackslash}p{2.9cm}>{\raggedright\arraybackslash}p{3.3cm}>{\raggedright\arraybackslash}p{2.0cm}>{\raggedright\arraybackslash}p{3.6cm}>{\raggedright\arraybackslash}p{2.1cm}@{}}
\toprule
& \textbf{Hypothèse} & \textbf{Statistique et seuil} & \textbf{Gel} & \textbf{Résultat} & \textbf{Verdict} \\
\midrule
N7 & Addendum : A, DiD IDE au sein des bassins de vie ; B, stabilité de N1 ; C, bande « pente » ;
D, cibles DREES &
A : seuils de N3, $Y \ge 5$\suppose{} ; B : même signe en 2016, 2018, 2021 ; D : résidu de Shapley
$\le 25$\suppose{}~\% &
04:37 décl. ; copie et résultats \textbf{dans le même commit} 6f5e034 &
A : T1 $p = \num{0.025}$\mesure, effet \num{-32.2}\mesure{}~\% ; C : couverture « 80~\% »
$\approx 96$\mesure{}~\% ; D2 : résidu 16\derive{}~\% &
A \textbf{échoué} ; B instable ; C CRPS passé, calibration \textbf{échouée} ; D1 \textbf{échoué},
D2 \textbf{passé} \\
\addlinespace
N8 & Bande $\log N \sim \mathcal N(\log \hat N, s^2 h)$ calibrée (forme suggérée par N7-C) &
$h = 5$\suppose, origines 2017--2021, $N \ge 30$\suppose{} ; couverture à 80~\% dans
$[75 ; 85]$\suppose{}~\% et à 50~\% dans $[44 ; 56]$\suppose{}~\% &
08:06 décl. ; copie et résultats \textbf{dans le même commit} 03a67d8 &
\num{1581}\mesure{} couples (SF, CD, PH, IDE) ; couverture \num{49.2}\mesure{}~\% et
\num{75.9}\mesure{}~\% ; $\Delta$CRPS $= \num{-0.0045}$\mesure &
\textbf{passé} dans l'échantillon, à \num{0.9}\derive{}~point de la borne ; \textbf{en attente des
données 2027} \\
\addlinespace
N9 & B2 : densité $\to$ installations des moins de 40~ans, $\beta < 0$ ; B1 : densité $\to$
rétention des 45--54~ans, $\beta > 0$ &
MCO à effets fixes, IC à 95~\% CR1 ; confirmé ou contraire selon l'IC entier ; placebo autres
professions &
08:10 décl. ; copie et résultats \textbf{dans le même commit} f1f4f41 &
B2 : IDE \num{-1.52} $[\num{-4.05} ; \num{+1.01}]$\mesure, MK et CD nuls\mesure{} ; B1 IDE
\num{+0.11} $[\num{+0.05} ; \num{+0.18}]$\mesure, placebo \num{-0.00}\mesure &
B2 \textbf{échoué} (signe non établi) ; B1 \textbf{passé} pour les IDE seuls, corrélationnel \\
\addlinespace
N10 & Au moins un grain infra-départemental a une bande P10--P90 calibrée &
couverture dans $[75 ; 85]$\suppose{}~\% dans chaque classe d'au moins 30\suppose{} unités &
08:31 décl. ; copie et résultats \textbf{dans le même commit} 167d15d &
aucun grain fiable ; couverture par classe 94\mesure{}~\% (IDE, 1--2 pros) à 22--59\mesure{}~\%
(plus de 100) ; K3 échoue aussi &
\textbf{échoué} \\
\addlinespace
N11 & K4 passe N10 puis l'origine réservée 2021, $h = 4$ & seuils de N10 ; les deux étapes &
08:50 décl. ; 08:50:22\mesure{} (a00cd67), avant le jugement (08:54\mesure) &
étape~1 : 2 à 6\mesure{} classes hors bande par grain ; étape~2 : 3 à 6\mesure &
\textbf{échoué} ; dépôt K4 \textbf{en attente des données 2027} \\
\addlinespace
N12 & K5 passe N10 puis l'origine réservée 2022, $h = 3$ & seuils de N10 ; dépôt écrit avant le
verdict &
09:05 décl. ; 09:05:46\mesure{} (2230ed0), \textbf{après} les soumissions (09:04:47 décl.) &
IDE au BV \num{65.9}\mesure{}~\%, EPCI \num{59.1}\mesure{}~\% ; 3 à 7\mesure{} classes hors bande &
\textbf{échoué} ; dépôt K5 \textbf{en attente des données 2027} \\
\bottomrule
\end{tabular}
\caption{Protocoles N7 à N12. Sources : \texttt{R/n7} l.21--220, \texttt{R/n8} l.28--70,
\texttt{R/n9} l.20--147, \texttt{R/n10} l.61--182, \texttt{R/n11} l.24--109, \texttt{R/n12}
l.19--83 ; heures de commit lues dans git.}
\label{val:tab-protocoles-b}
\end{table}

\paragraph{Bilan.} Aucun protocole ne passe sans réserve. Les seules affirmations admises sont
étroites : pour la variante par défaut, critère N1 passé au poolé, victoire portée par l'origine
2021 seule (\texttt{D} l.1828--1830) ; le moteur respecte ses cas limites analytiques
(\texttt{R/n2} l.165--188) ; la prédiction marche/arrêt des kinés tient au département
(\texttt{R/n3} l.138--146) ; la bande en $\sqrt h$ est calibrée dans l'échantillon pour SF, CD, PH
et IDE (\texttt{R/n8} l.45--70), mais le code servi ne l'affiche pas (section~\ref{val:perimetre}). Ces verdicts sont rapportés tels que les protocoles les donnent.
L'un d'eux n'est pas qualifié ainsi dans les rapports : l'échec du critère de bande de N4, qui
est rapporté (14\mesure{} classes sur 21\mesure) sans le mot « échoue » (annexe~\ref{val:ann}).

\subsection{Intégrité du pré-enregistrement}\label{val:integrite}

Le pré-enregistrement ne vaut que si l'on peut prouver que le gel précède le calcul. Dans ce
dépôt, la seule preuve disponible est l'historique git. Elle est insuffisante pour la moitié des
protocoles.

\paragraph{Six protocoles sur douze n'ont pas d'antériorité prouvée par git.}
Sur 12\mesure{} protocoles, 6\mesure{} ont un gel dont git ne montre pas qu'il précède le calcul
(lecture de \texttt{git log -- v2/validation/preregistered}, 2026-11-15).
\begin{itemize}
  \item \textbf{N3.} La copie gelée et \texttt{out/did\_results.json} entrent dans le même commit
        316a93d, à 03:25:50\mesure. Le rapport date le calcul vers « 03:0x » (déclaré), donc avant
        le premier commit. Le gel de 02:54 repose sur des dates de fichiers et sur une
        auto-déclaration. La copie gelée porte de plus l'en-tête « lancé pour les IDE ».
  \item \textbf{N7.} Copie et résultats dans le même commit 6f5e034, à 04:46:59\mesure. Le gel
        déclaré (04:37) est postérieur au résultat de N3, ce qui est assumé : N7 est un addendum.
  \item \textbf{N8.} Copie et résultats dans le même commit 03a67d8, à 08:08:49\mesure.
  \item \textbf{N9.} Copie et résultats dans le même commit f1f4f41, à 08:12:33\mesure.
  \item \textbf{N10.} Copie et résultats dans le même commit 167d15d, à 08:33:07\mesure.
  \item \textbf{N12.} La copie est commitée seule (2230ed0, 09:05:46\mesure), avant le jugement.
        Mais les soumissions de K5, datées 09:04:47 et 09:05:00 (déclaré), précèdent ce commit. Le
        protocole n'est donc pas prouvé antérieur aux soumissions qu'il juge.
\end{itemize}
Pour les six autres (N1, N2, N4, N5, N6, N11), la copie gelée est commitée avant le commit des
résultats. C'est une preuve d'ordre relative, à la réserve suivante près.

\paragraph{L'historique a été réécrit.} Le 2026-11-14, cinq commits ont été réécrits avec
l'accord de l'utilisatrice ou de l'utilisateur, par \texttt{git commit-tree} et
\texttt{git update-ref}, pour en retirer une ligne d'attribution, à arbres identiques
(\texttt{D} l.1158--1164). L'un d'eux, b2756d8 devenu 72812c8, est le commit qui gèle N5 et N6.
Ses dates d'auteur et de commit ont été conservées (03:55:53\mesure{} toutes deux, lues par
\texttt{git log}). Cela montre que les dates git se fixent librement : un historique réécrit peut
porter n'importe quelle date. Aucun horodatage externe n'est mentionné dans les fichiers lus. La
preuve d'ordre de git est donc une preuve de bonne foi, pas une preuve opposable. Par ailleurs,
la copie gelée du plan a une date de modification à 10:45 (lue par \texttt{stat} sur le fichier) ; son empreinte
est intacte.

\paragraph{Les empreintes concordent toutes.} Le 2026-11-15, les 13\mesure{} empreintes SHA-256
des copies gelées ont été recalculées (\texttt{shasum -a 256 v2/validation/preregistered/*.md}).
Les 12\suppose{} premiers caractères égalent ceux qui sont inscrits, pour le plan
(\texttt{0108455d34aa}) et pour N1 à N12. Les empreintes des dépôts de prévision concordent aussi :
\texttt{ca036374d216} (variante par défaut), \texttt{039d21878b7a} (K4), \texttt{0169b97d4fac} (K5),
\texttt{4f7400845340} (A0 de l'axe~2), \texttt{689b230b5b50} (C1 de l'axe~2, recalculée le
2026-11-15). Ce contrôle prouve que les copies n'ont pas changé depuis
l'inscription de l'empreinte. Il ne prouve pas que l'empreinte a été inscrite avant le calcul :
elle l'est dans des fichiers du même dépôt.

\paragraph{Mesures proposées pour les prochains dépôts.} Les prix sont proposés ici ; ils sont
supposés\suppose.
\begin{enumerate}
  \item \textbf{Horodatage externe.} Soumettre l'empreinte de chaque copie gelée à un service
        d'horodatage indépendant (jeton RFC~3161 d'une autorité d'horodatage, ou ancrage
        OpenTimestamps) au moment du gel, et commiter le jeton. Prix : $\sim 10$\suppose{}~min
        Claude par lot, plus la création d'un compte si le service l'exige (décision humaine).
  \item \textbf{Pousser le gel seul.} Commiter et pousser la copie gelée sur le dépôt distant,
        dans un commit qui ne contient rien d'autre, avant de lancer le calcul. L'heure de
        réception par l'hébergeur sert de second témoin. Prix : nul en temps Claude ; exige une
        autorisation de push à chaque gel.
  \item \textbf{Registre tiers.} Pour un protocole destiné à publication, déposer la copie sur un
        registre de pré-enregistrement public (par exemple OSF Registries), qui date le dépôt
        hors du contrôle de l'équipe. Prix : temps humain, quelques minutes par protocole.
  \item \textbf{Ordre vérifiable par script.} Ajouter un contrôle qui refuse le commit d'un
        fichier de résultats si la copie gelée correspondante n'est pas déjà dans un commit
        antérieur. Prix : $\sim 20$\suppose{}~min Claude.
  \item \textbf{Ne jamais réécrire} un commit qui contient une copie gelée, ou consigner
        l'ancienne et la nouvelle empreinte de commit dans le registre externe.
\end{enumerate}

\subsection{Ce qui est validé, ce qui ne l'est pas}\label{val:perimetre}

\paragraph{Prévision départementale à 5~ans (N1 ; $h = 1$ à 10 en secondaire).} Critère N1 passé au
poolé : $\Delta = \num{-0.63}$ $[\num{-0.97} ; \num{-0.27}]$\mesure{} au périmètre privé (\texttt{R/n1}
l.160--166) ; victoire portée par l'origine 2021 seule. En 2016 et 2018, l'IC contient 0\mesure, et le
$\tau$ de Kendall entre classements d'origines vaut au plus \num{0.29}\mesure{} (\texttt{R/n7}
l.118--220). Seul $h = 5$ est jugé ; aux horizons secondaires, la persistance fait mieux qu'A2 à $h = 8$
et 10 (\num{13.56} et \num{16.06}\mesure{}~\% contre \num{14.33} et \num{17.59}\mesure{}~\%,
\texttt{v2/methodes/out/c1\_gbm\_verdicts.json}, sans IC). Le garde-fou à $h = 1$ est rapporté :
\texttt{limite\_h1\_ok} vaut \texttt{true} dans \texttt{v2/methodes/out/verdicts\_prive.json} et
\texttt{coh\_verdicts.json}.

\paragraph{Le modèle servi n'est pas celui que l'axe~3 a jugé.} Les protocoles N1 à N8 jugent A0, la
variante par défaut. Le jumeau sert C1, l'estimateur de cohortes de l'axe~2 (section~\ref{met:bayes}).
C1 passe le critère de N1 au banc du 2026-11-14 : $\Delta = \num{-1.57}$ $[\num{-2.17} ;
\num{-0.99}]$\mesure{} contre la tendance (\texttt{v2/methodes/out/coh\_verdicts.json}). Deux de ses
cellules (CD et SF, origine 2021) ont été vues par son pilote avant la correction de la forme des
variances. Il échoue au critère de couverture de N6 : sa bande à 80~\% couvre \num{73.5}\mesure{}~\%
pour une cible de $[75 ; 85]$\suppose{}~\% (\texttt{calibre80} = \texttt{false}), et ce score n'est pas
hors échantillon (deux des trois cellules jugées sont celles du pilote). Son dépôt 2027 est
\texttt{v2/\allowbreak validation/\allowbreak depot/\allowbreak axe2\_cohortes\_c1\_2026-11-14.parquet}, empreinte
\texttt{689b230b5b50} (section~\ref{val:depot}).

\paragraph{Le grain sous le département n'est pas validé.} Aucune variante communale n'a de bande
calibrée à la commune, au bassin de vie ou à l'EPCI : ni les estimateurs de référence ni K3
(N10), ni K4 (N11), ni K5 (N12) (\texttt{R/n10} l.84--182, \texttt{R/n11} l.57--109,
\texttt{R/n12} l.43--83). Le seuil $N_{\min} = 51$\mesure{} de N4 est obtenu sous un oracle
départemental qui connaît la vraie valeur de la cible ; la version à projection réelle n'est pas
lancée (\texttt{O} l.138). En conséquence, l'affichage retenu est le département ; à la commune,
seul l'observé s'affiche (\texttt{D} l.2151--2194). Les sages-femmes et les pharmaciens n'ont pas
de série communale.

\paragraph{Aucun levier n'est validé de bout en bout.} Le levier de régulation infirmière échoue
deux fois au test de tendances parallèles fixé d'avance : $p = \num{0.043}$\mesure{} au
département, $p = \num{0.025}$\mesure{} au bassin de vie, pour un seuil de \num{0.10}\suppose{}
(\texttt{R/n3} l.165--224, \texttt{R/n7} l.123--153). L'effet estimé est grand
(\num{-30.3}\mesure{} log-points), mais le design ne permet pas de l'attribuer. Seule la prédiction
marche/arrêt des kinés passe, au département (écart des pentes \num{+3.1}\mesure{} pts/an), après
l'annulation par le Conseil d'État \citep{ce2014mk}. Au bassin de vie, son test de tendances
parallèles échoue ($p = \num{0.059}$\mesure). Le résultat IRDES sur les infirmières
\citep{irdes270} reste une source externe, pas une validation du module.

\paragraph{$\theta = \num{0.80}$ est utilisé par défaut malgré un verdict « non interprétable ».}
N3 n'autorise l'interprétation de $\theta = \beta_{\text{obs}}/\beta_{\text{mod}}$ que si T1 à T4
passent ; T1 échoue, donc $\theta = \num{0.80}$ $[\num{0.74} ; \num{0.87}]$\mesure{} n'est pas
interprété (\texttt{R/n3} l.189--191). Ce nombre est pourtant devenu le paramètre par défaut
\texttt{theta} du module de régulation, étiqueté « supposé » (\texttt{D} l.1189--1192), avec une
plage $[\num{0.5} ; 1]$\suppose{} (\texttt{D} l.1064). L'étiquette est honnête ; l'usage dépasse
ce que le protocole autorise. Un utilisateur qui lit « \num{0.80} » doit savoir que ce nombre sort
d'un design invalide. Deux faits limitent la portée : au département et au-dessus, la régulation
par défaut n'a aucun effet, car les entrants bloqués restent dans la cellule (\texttt{D}
l.1198--1200) ; à la commune, $\theta$ agit pleinement.

\paragraph{2050 n'est pas validé.} Les cas limites analytiques L1 à L3 passent (\texttt{R/n2}
l.165--188). Mais la reproduction du tendanciel DREES \citep{drees2021dd76}, jouée en forme
affaiblie faute de série de quotas, échoue : \num{-14.1}\mesure, \num{-15.0}\mesure{} et
\num{-17.4}\mesure{}~\% en 2050 pour CD, PH et SF (\texttt{R/n2} l.133--163). La comparaison à la
projection infirmière \citep{er1319} échoue aussi (\num{-25.2}\mesure{}~\% en 2050). La forme
fidèle de N2-A n'est pas jouée. L'hypothèse d'entrées domine l'écart : une variante « filière »,
exploratoire, se place \num{+14} à \num{+15}\mesure{}~\% au-dessus de DD76. Selon la règle de N2,
un échec bloque tout chiffre 2050 pour la profession : aucun chiffre 2050 n'est citable
(\texttt{D} l.1841--1846).

\paragraph{La bande affichée n'est pas celle que N8 calibre.} Le code servi n'affiche pas la bande en
$\sqrt h$ de N8 : il affiche $v\,(1 \pm k_p h)$, où $k_p$ est la pente de l'erreur du total national de
C1 (section~\ref{met:servi} ; \texttt{v2/prototype/server.py} l.158--180,
\texttt{v2/prototype/index.html} l.225). La bande calibrée par N8 n'est donc affichée pour aucune
profession. De plus, N8 calibre la bande en $\sqrt h$ pour SF, CD, PH et IDE seulement : les kinés ne figurent pas dans sa liste de professions (\texttt{R/n8} l.8
et l.48). Le rapport ne donne pas la raison ; une raison plausible est qu'à $h = 5$ les kinés n'ont
qu'une\derive{} origine jugeable, leur série RPPS commençant en 2017 (\texttt{A} l.127). Les bandes de N6 (trop étroites) et de N7-C (trop
larges, $\approx 96$\mesure{}~\% pour « 80~\% ») échouent. Le rapport N10 (l.92) et l'article
(\texttt{A} l.599--600) disent pourtant que l'affichage départemental porte « la bande $\sqrt h$
calibrée de N8 » : c'est inexact pour toutes les professions, et pour les kinés aucune bande
départementale n'est calibrée.

\paragraph{Le moteur sous bruit.} Le moteur est exact sans bruit (M1, écart maximal
$4\cdot10^{-14}$\mesure{}~\%). Sous bruit de Poisson, il est biaisé : M3 a une MAPE départementale
de \num{10.2}\mesure{}~\% pour un plancher de \num{5.7}\derive{}~\%, et la SBC des entrées échoue
($p = 3\cdot10^{-46}$\mesure). Le mécanisme proposé est le plafonnement \texttt{min(r, 1)} des taux
de rétention, l'excédent étant compté en entrées (\texttt{D} l.1847--1849).

\paragraph{Boucles.} B1 (densité $\to$ rétention) n'est confirmé que pour les IDE libéraux, de façon
corrélationnelle : \num{+0.11}\mesure{} sur l'échelle du protocole, soit $\approx +8$\derive{}~\% de
rétention quand la densité double. B2 n'a de signe établi pour aucune profession
(\texttt{R/n9} l.109--147). Rien n'est causal.

\subsection{Le dépôt de prévision 2027}\label{val:depot}

\paragraph{Ce qui est déposé.} Le 2026-11-14, la variante par défaut a déposé ses prévisions de
l'origine 2026 pour 2027, 2028 et 2030, par profession et département, au périmètre privé, avec
les quantiles 10 et 90~\% de la bande en $\sqrt h$ :
\texttt{v2/\allowbreak validation/\allowbreak depot/\allowbreak jumeau\_taux\_locaux\_2026-11-14.parquet}, empreinte
\texttt{ca036374d216} (\texttt{D} l.1964--1978). Totaux nationaux 2027 : CD
\num{47835}\mesure, MK \num{104755}\mesure, PH \num{63311}\mesure, SF \num{13301}\mesure{} (sorties
du modèle). Les infirmiers libéraux en sont exclus, leur série RPPS s'arrêtant en 2023.
Le même jour, C1, le modèle servi, a déposé
\texttt{v2/\allowbreak validation/\allowbreak depot/\allowbreak axe2\_cohortes\_c1\_2026-11-14.parquet} (empreinte
\texttt{689b230b5b50}, commit 470acf1) : les cinq professions, avec les quantiles 10, 50 et 90~\%.
Les IDE y partent de l'origine 2023 (cibles 2027, 2028 et 2030, soit $h = 4$, 5 et 7\derive) ; leur
score suppose une publication RPPS des IDE libéraux après 2023, qui n'est pas garantie. Les
dépôts K4 (\texttt{039d21878b7a}) et K5 (\texttt{0169b97d4fac}) portent sur la commune. Le score
sera calculé à la publication du RPPS 2027, attendue vers juillet 2027 \citep{drees2026demo}, sans
toucher aux fichiers. C'est le seul test qui ne soit contaminé ni par K1 ni par le choix de la
bande en $\sqrt h$.

\paragraph{Limites.}
\begin{itemize}
  \item \textbf{Il ne jouera pas le critère principal.} Depuis l'origine 2026, les cibles 2027,
        2028 et 2030 donnent $h = 1$, 2 et 4\derive. Le critère de N1 est à $h = 5$\suppose. Le
        protocole N1 annonce pourtant un score « avec le même critère » (\texttt{R/n1}
        l.115--116), ce qui est impossible. Il faudra déclarer, avant 2027, quel critère s'applique
        à $h \le 4$.
  \item \textbf{Sa composition est inversée par rapport à la calibration.} Il inclut les kinés,
        dont la bande n'a jamais été calibrée, et exclut les IDE, que N8 a calibrés.
  \item \textbf{Une seule origine.} Le score portera sur une origine et au plus trois horizons :
        il dira si 2026 est une année « facile » ou non, pas si la méthode est stable.
  \item \textbf{K4 et K5 ne peuvent plus être confirmés.} Leurs dépôts partent de l'origine 2025
        et non 2026 (écart d'exécution déclaré). Ils mesureront une couverture, mais N11 et N12
        ont déjà rendu un verdict d'échec.
  \item \textbf{Calendrier.} La date de publication de CartoSanté 2027 est inconnue ; le prix du
        score est $\sim 15$\suppose{}~min Claude, mais l'attente en temps du monde va au moins
        jusqu'en juillet 2027 (\texttt{A} l.744).
  \item \textbf{Nom sans préfixe.} Le fichier du dépôt garde un nom sans préfixe d'axe, contrairement
        à la règle de \texttt{D} l.2377, et ce nom a déjà été écrasé une fois (\texttt{D}
        l.2361--2373) ; son empreinte actuelle concorde avec celle inscrite.
\end{itemize}

\subsection{Menaces à la validité}\label{val:menaces}

\paragraph{Antériorité des gels.} Voir la section~\ref{val:integrite} : 6\mesure{} protocoles sur
12\mesure{} sans preuve git d'antériorité, historique réécrit, aucun horodatage externe.

\paragraph{Fuite temporelle.}
\begin{itemize}
  \item N1 : les cibles 2023 et 2026 ont été vues par K1 ; le banc rétrospectif est en partie dans
        l'échantillon de conception (\texttt{R/n1} l.137--140).
  \item N8 : la forme en $\sqrt h$ a été suggérée par N7-C et testée sur les mêmes origines.
  \item K3 a été inscrit à 08:41 (déclaré), après les résultats de référence de N10 (08:33\mesure),
        et sa conception cite ces résultats (\texttt{D} l.3007--3009). K4 et K5 ont été conçus
        après les jugements précédents : risque de sur-ajustement au juge.
  \item Origines réservées : l'origine 2021 partage 4\derive{} années sur 5 avec une fenêtre vue,
        l'origine 2022 3\derive{} sur 3. Ce sont des confirmations partielles, sur des périodes
        où la tendance calibre déjà 8\mesure{} classes sur 8.
  \item N2-A, variante filière : elle utilise des inscrits connus après 2020.
  \item N4 : l'oracle départemental retire $\approx 8$\mesure{}~points de MAPE (K1) ; il sous-estime
        l'erreur communale.
\end{itemize}

\paragraph{Cellules vues avant inscription.} N10 a été gelé après le passage de N4 sur les mêmes
données communales, pour des origines qui se recouvrent (2014--2020) : l'échec de la bande dans les
grandes classes était connu. N7-A a été écrit après le résultat de N3. La variante \texttt{vivier}
a été évaluée sur une seule origine, hors pré-enregistrement (\texttt{D} l.1707--1711).

\paragraph{Multiplicité.} 12\mesure{} protocoles et plusieurs dizaines de critères. Bonferroni est
appliqué à l'intérieur de certains bancs : la SBC, les cinq variantes stock-flux (IC à 99~\%) et le
GBM et le MLP (IC à \num{97.5}~\%). Il ne l'est pas entre protocoles. Le registre N1 compte au moins 8\mesure{}
estimateurs. Le levier IDE a été testé 2\mesure{} fois, les kinés 2\mesure{} fois, et cinq
variantes communales (K à K5) ont été soumises. N9 fait 8\mesure{} tests de signe à 95~\%, soit
$\approx \num{0.4}$\derive{} faux positif attendu sous $H_0$ ($8 \times \num{0.05}$) : le seul
« confirmé », B1 pour les IDE, se lit avec cette réserve.

\paragraph{Petits échantillons.} E1 repose sur $G = 59$\mesure{} départements, E2 sur 26 à
29\mesure, B1 des sages-femmes sur 15\mesure. N6 n'est jugeable que sur une origine, trois
professions et 279\derive{} couples. À $h = 5$, N1 dispose de 6\derive{} origines pour SF, CD et
PH, 3\derive{} pour les IDE et 1\derive{} pour les kinés. N5 utilise 20\suppose{} réplications.

\paragraph{Intervalles probablement trop étroits.} Le bootstrap de N1 et N6 tire les mêmes
départements pour toutes les cellules (\texttt{v2/validation/forecast.py} l.58--65), ce qui conserve
la corrélation entre origines d'un même département. Il ne rééchantillonne pas les origines (6\mesure{}
au plus), ignore le choc commun à tous les départements d'une origine et la corrélation entre
départements (\texttt{R/n1} l.141--142). Le remède est de rééchantillonner aussi les origines, ou un
bootstrap à deux grappes (départements et origines), et non des blocs département $\times$ origine.
Une correction du bootstrap est inscrite au registre (\texttt{O} l.67, $\sim 20$\suppose{}~min).

\paragraph{Critères au bord du seuil.} N8 passe à \num{75.9}\mesure{}~\% pour une borne de
75\suppose{}~\%. Le test T4 au bassin de vie passe « de justesse » (\texttt{R/n7}).

\paragraph{Écarts d'exécution déclarés.} N2-A en forme affaiblie ; N6 avec 100\mesure{} tirages
au lieu de 200\suppose{} pour l'APC ; dépôts K4 et K5 à l'origine 2025 ; K5 avec un $\sigma_D$
robuste, changé avant passage ; N5 avec deux précisions sur M4 fixées sans lecture des résultats
(déclaré).

\paragraph{Dérive des données après le gel.} La table privée des sages-femmes a été modifiée après
N1 : la MAPE de A2 passe de \num{8.57} à \num{8.60}\mesure{} (\texttt{O} l.195). Les chiffres de N1
ne se reproduisent plus exactement sur les données actuelles ; une relance coûte
$\sim 10$\suppose{}~min machine.

\paragraph{Identification des leviers.} La date du zonage Atlasanté \citep{atlasante2014zonage}
n'est que « probablement 2012 », par déduction ; si elle est fausse, E1 et E2 tombent ensemble.
Les dates d'effet viennent de sources secondaires. « Pas de régulation dentaire avant 2025 » est
supposé. Le retour à la moyenne peut expliquer jusqu'à $\sim 1/7$\derive{} de l'effet IDE. Au
bassin de vie, le débordement vers les communes voisines fait de $\beta$ une borne haute de
l'effet sur le total. Une borne de sensibilité à la Rambachan--Roth \citep{rambachan2023} est
rapportée dans l'article (\texttt{A} l.402).

\paragraph{Différences de champ.} DD76 porte sur tous les modes d'exercice, le jumeau sur le
privé ; ER~1319 sur les infirmières en emploi, le jumeau sur le RPPS libéral ; B2 sur CartoSanté,
B1 sur le RPPS libéral et mixte.

\paragraph{Seuils choisis, non dérivés.} 3\suppose{}~\% (N2), 25\suppose{}~\% (N4),
$p \ge \num{0.10}$\suppose{} (N3), $[75 ; 85]$\suppose{}~\% et $N \ge 30$\suppose{} sont des
choix. Un autre seuil changerait certains verdicts, en particulier N8.
\endgroup
